% Generated draft: requires mathematical and source review.
\documentclass[11pt]{article}
\usepackage[margin=25mm]{geometry}
\usepackage{fontspec}
\setmainfont{Latin Modern Roman}
\usepackage{amsmath,amssymb,mathtools,unicode-math}
\setmathfont{Latin Modern Math}
\usepackage{enumitem,needspace,xurl,hyperref}
\setlength{\parindent}{0pt}
\setlength{\parskip}{5pt}
\setlength{\emergencystretch}{3em}
\newcommand{\N}{\mathbb{N}}
\newcommand{\Z}{\mathbb{Z}}
\newcommand{\Q}{\mathbb{Q}}
\newcommand{\R}{\mathbb{R}}
\newcommand{\C}{\mathbb{C}}
\newcommand{\F}{\mathbb{F}}
\newcommand{\A}{\mathbb{A}}
\newcommand{\PP}{\mathbb{P}}
\newcommand{\E}{\mathbb{E}}
\newcommand{\eps}{\varepsilon}
\newcommand{\dd}{\,\mathrm{d}}
\begin{document}
\section*{UnsolvedMath ID: 30006653}
\textbf{Source problem code:} catalogue entry
\subsection{Definitions and mathematical statement}
Let \(p\) be a prime and let \(G=(\Z/p\Z)^r\) be an elementary abelian \(p\)-group of rank \(r\in\Z_{\geq 0}\). A continuous action of \(G\) on a space \(X\) is free if \(gx=x\) implies \(g=e\) for every \(g\in G\) and \(x\in X\). A finite-dimensional CW complex has a finite bound on its cell dimensions; this does not require finitely many cells. Write \(X\simeq Y\) for ordinary homotopy equivalence.

\textbf{Rank-bound assertion, with the sphere-dimension convention unresolved.} In the positive-dimensional, nonempty-product interpretation, the assertion is: for every prime \(p\), every \(r\in\Z_{\geq 0}\), every integer \(k\geq 1\), every collection of integers \(n_1,\ldots,n_k\geq 1\), and every finite-dimensional CW complex \(X\) admitting a free continuous \(G\)-action, the implication
\[
X\simeq\prod_{i=1}^{k}S^{n_i}
\quad\Longrightarrow\quad r\leq k
\]
holds. Here \(S^{n_i}\) is the \(n_i\)-dimensional sphere. The supplied record does not explicitly specify whether zero-dimensional spheres or the empty product are included, so this displayed formulation states an explicit interpretation rather than resolving those boundary conventions. The record also does not specify whether additional regularity, such as cellularity of the action, is intended.

\textbf{Highlighted equidimensional targets.} For each \(n\in\{3,7\}\), each integer \(k\geq 1\), and each \(r\in\Z_{\geq 0}\), establish that a free action of \((\Z/2\Z)^r\) on a finite-dimensional CW complex satisfying
\[
X\simeq(S^n)^k
\]
implies \(r\leq k\).

\textbf{Scope and status qualifications.} The supplied reconstructed statement allows sphere dimensions to vary, whereas the supplied original excerpt explicitly discusses only \(X\simeq(S^n)^k\). The excerpt reports that the equidimensional conjecture is settled except when \(p=2\) and \(n\in\{3,7\}\). The supplied dated assessment describes the general conjecture as partially solved, including established odd-prime equidimensional cases, with full generality and the highlighted cases remaining open. These are supplied source assessments, not independently verified mathematical or literature conclusions.
\subsection{Short English statement}
The proposed bound says that a freely acting elementary abelian \(p\)-group has rank at most the number of sphere factors in the homotopy type of the space. The record particularly asks for the cases \(X\simeq(S^3)^k\) and \(X\simeq(S^7)^k\) when \(p=2\). The broader formulation, boundary conventions, and intended action hypotheses require human review.
\subsection{Sources}
\begin{itemize}
\item Local UnsolvedMath problem registry: supplied complete record for permanent ID \(30006653\), source problem code \texttt{catalogue entry}, titled \emph{Rank Conjecture for Free Elementary-Abelian Actions}, in the Oberwolfach Reports Open Problems collection. No individual proposer is supplied; the recorded proposed year is 2026.
\item Originating report: \emph{Cohomology of Finite Groups: Interactions and Applications} (2026), Oberwolfach Reports, pp.~541--610, DOI \texttt{10.4171/owr/2026/9}. The supplied original excerpt identifies p.~567. \url{https://doi.org/10.4171/owr/2026/9}.
\item Supplied literature reference: Carlsson, \emph{On the Rank of Abelian Groups Acting Freely on Products of Spheres}. \url{https://eudml.org/doc/142960}.
\item Supplied literature reference: Adem--Browder, \emph{Actions of elementary Abelian p-groups}. \url{https://doi.org/10.1016/0040-9383(88)90024-9}.
\item Literature assessment checked at \texttt{2026-08-22}: the supplied curation labels the problem partially solved and reports the rank bound established in broad cases, including all odd primes in the cited equidimensional setting, while full generality and the cases \(p=2\), \(n=3,7\) remain open. This is supplied dated curation derived from the Oberwolfach record, not an independently verified assessment; time-sensitive status and cited claims require independent verification.
\end{itemize}
\end{document}
